% Options for packages loaded elsewhere
\PassOptionsToPackage{unicode}{hyperref}
\PassOptionsToPackage{hyphens}{url}
%
\documentclass[
]{article}
\usepackage{amsmath,amssymb}
\usepackage{iftex}
\ifPDFTeX
  \usepackage[T1]{fontenc}
  \usepackage[utf8]{inputenc}
  \usepackage{textcomp} % provide euro and other symbols
\else % if luatex or xetex
  \usepackage{unicode-math} % this also loads fontspec
  \defaultfontfeatures{Scale=MatchLowercase}
  \defaultfontfeatures[\rmfamily]{Ligatures=TeX,Scale=1}
\fi
\usepackage{lmodern}
\ifPDFTeX\else
  % xetex/luatex font selection
\fi
% Use upquote if available, for straight quotes in verbatim environments
\IfFileExists{upquote.sty}{\usepackage{upquote}}{}
\IfFileExists{microtype.sty}{% use microtype if available
  \usepackage[]{microtype}
  \UseMicrotypeSet[protrusion]{basicmath} % disable protrusion for tt fonts
}{}
\makeatletter
\@ifundefined{KOMAClassName}{% if non-KOMA class
  \IfFileExists{parskip.sty}{%
    \usepackage{parskip}
  }{% else
    \setlength{\parindent}{0pt}
    \setlength{\parskip}{6pt plus 2pt minus 1pt}}
}{% if KOMA class
  \KOMAoptions{parskip=half}}
\makeatother
\usepackage{xcolor}
\usepackage[margin=1in]{geometry}
\usepackage{color}
\usepackage{fancyvrb}
\newcommand{\VerbBar}{|}
\newcommand{\VERB}{\Verb[commandchars=\\\{\}]}
\DefineVerbatimEnvironment{Highlighting}{Verbatim}{commandchars=\\\{\}}
% Add ',fontsize=\small' for more characters per line
\usepackage{framed}
\definecolor{shadecolor}{RGB}{248,248,248}
\newenvironment{Shaded}{\begin{snugshade}}{\end{snugshade}}
\newcommand{\AlertTok}[1]{\textcolor[rgb]{0.94,0.16,0.16}{#1}}
\newcommand{\AnnotationTok}[1]{\textcolor[rgb]{0.56,0.35,0.01}{\textbf{\textit{#1}}}}
\newcommand{\AttributeTok}[1]{\textcolor[rgb]{0.13,0.29,0.53}{#1}}
\newcommand{\BaseNTok}[1]{\textcolor[rgb]{0.00,0.00,0.81}{#1}}
\newcommand{\BuiltInTok}[1]{#1}
\newcommand{\CharTok}[1]{\textcolor[rgb]{0.31,0.60,0.02}{#1}}
\newcommand{\CommentTok}[1]{\textcolor[rgb]{0.56,0.35,0.01}{\textit{#1}}}
\newcommand{\CommentVarTok}[1]{\textcolor[rgb]{0.56,0.35,0.01}{\textbf{\textit{#1}}}}
\newcommand{\ConstantTok}[1]{\textcolor[rgb]{0.56,0.35,0.01}{#1}}
\newcommand{\ControlFlowTok}[1]{\textcolor[rgb]{0.13,0.29,0.53}{\textbf{#1}}}
\newcommand{\DataTypeTok}[1]{\textcolor[rgb]{0.13,0.29,0.53}{#1}}
\newcommand{\DecValTok}[1]{\textcolor[rgb]{0.00,0.00,0.81}{#1}}
\newcommand{\DocumentationTok}[1]{\textcolor[rgb]{0.56,0.35,0.01}{\textbf{\textit{#1}}}}
\newcommand{\ErrorTok}[1]{\textcolor[rgb]{0.64,0.00,0.00}{\textbf{#1}}}
\newcommand{\ExtensionTok}[1]{#1}
\newcommand{\FloatTok}[1]{\textcolor[rgb]{0.00,0.00,0.81}{#1}}
\newcommand{\FunctionTok}[1]{\textcolor[rgb]{0.13,0.29,0.53}{\textbf{#1}}}
\newcommand{\ImportTok}[1]{#1}
\newcommand{\InformationTok}[1]{\textcolor[rgb]{0.56,0.35,0.01}{\textbf{\textit{#1}}}}
\newcommand{\KeywordTok}[1]{\textcolor[rgb]{0.13,0.29,0.53}{\textbf{#1}}}
\newcommand{\NormalTok}[1]{#1}
\newcommand{\OperatorTok}[1]{\textcolor[rgb]{0.81,0.36,0.00}{\textbf{#1}}}
\newcommand{\OtherTok}[1]{\textcolor[rgb]{0.56,0.35,0.01}{#1}}
\newcommand{\PreprocessorTok}[1]{\textcolor[rgb]{0.56,0.35,0.01}{\textit{#1}}}
\newcommand{\RegionMarkerTok}[1]{#1}
\newcommand{\SpecialCharTok}[1]{\textcolor[rgb]{0.81,0.36,0.00}{\textbf{#1}}}
\newcommand{\SpecialStringTok}[1]{\textcolor[rgb]{0.31,0.60,0.02}{#1}}
\newcommand{\StringTok}[1]{\textcolor[rgb]{0.31,0.60,0.02}{#1}}
\newcommand{\VariableTok}[1]{\textcolor[rgb]{0.00,0.00,0.00}{#1}}
\newcommand{\VerbatimStringTok}[1]{\textcolor[rgb]{0.31,0.60,0.02}{#1}}
\newcommand{\WarningTok}[1]{\textcolor[rgb]{0.56,0.35,0.01}{\textbf{\textit{#1}}}}
\usepackage{graphicx}
\makeatletter
\def\maxwidth{\ifdim\Gin@nat@width>\linewidth\linewidth\else\Gin@nat@width\fi}
\def\maxheight{\ifdim\Gin@nat@height>\textheight\textheight\else\Gin@nat@height\fi}
\makeatother
% Scale images if necessary, so that they will not overflow the page
% margins by default, and it is still possible to overwrite the defaults
% using explicit options in \includegraphics[width, height, ...]{}
\setkeys{Gin}{width=\maxwidth,height=\maxheight,keepaspectratio}
% Set default figure placement to htbp
\makeatletter
\def\fps@figure{htbp}
\makeatother
\setlength{\emergencystretch}{3em} % prevent overfull lines
\providecommand{\tightlist}{%
  \setlength{\itemsep}{0pt}\setlength{\parskip}{0pt}}
\setcounter{secnumdepth}{-\maxdimen} % remove section numbering
\ifLuaTeX
  \usepackage{selnolig}  % disable illegal ligatures
\fi
\usepackage{bookmark}
\IfFileExists{xurl.sty}{\usepackage{xurl}}{} % add URL line breaks if available
\urlstyle{same}
\hypersetup{
  hidelinks,
  pdfcreator={LaTeX via pandoc}}

\author{}
\date{\vspace{-2.5em}}

\begin{document}

\section{Enchainement Limit Orders
(EURGBP)}\label{enchainement-limit-orders-eurgbp}

\subsection{Vue rapide}\label{vue-rapide}

\begin{enumerate}
\def\labelenumi{\arabic{enumi}.}
\tightlist
\item
  Une equipe envoie un ordre (\texttt{POST\ /trade/EURGBP}) avec
  \texttt{price} =\textgreater{} ordre \textbf{LIMIT}.
\item
  \texttt{Exchange.try\_trade(...)} met a jour la liquidite NPC puis
  tente le matching.
\item
  Si le prix limite croise deja le carnet, l'ordre est execute
  (total/partiel).
\item
  Sinon, le reliquat est place en \textbf{resting order} dans l'order
  book.
\item
  Quand une autre equipe envoie un \textbf{market order}, elle consomme
  les meilleures offres.
\item
  Les deux cotes du trade sont comptabilises (agresseur + contrepartie
  passive).
\item
  Le GUI lit \texttt{/orderbook} et \texttt{/tradeHistory} pour afficher
  carnet + historique.
\end{enumerate}

\subsection{Diagramme Sequence
(Mermaid)}\label{diagramme-sequence-mermaid}

\begin{Shaded}
\begin{Highlighting}[]
\NormalTok{sequenceDiagram}
\NormalTok{    autonumber}
\NormalTok{    participant TA as Team A (Seller)}
\NormalTok{    participant TB as Team B (Buyer)}
\NormalTok{    participant API as TradeEndpoint}
\NormalTok{    participant EX as Exchange}
\NormalTok{    participant NPC as NPCManager}
\NormalTok{    participant OB as OrderBook}
\NormalTok{    participant H as TradeHistory}

\NormalTok{    TA{-}\textgreater{}\textgreater{}API: POST /trade/EURGBP \{side:sell, qty, price\}}
\NormalTok{    API{-}\textgreater{}\textgreater{}EX: try\_trade(trader\_id, side, qty, limit\_price)}
\NormalTok{    EX{-}\textgreater{}\textgreater{}NPC: update(current\_price)}
\NormalTok{    EX{-}\textgreater{}\textgreater{}OB: resolve\_crossed\_book()}
\NormalTok{    EX{-}\textgreater{}\textgreater{}OB: execute\_limit(...)}

\NormalTok{    alt Limit croise deja le carnet}
\NormalTok{        OB{-}{-}\textgreater{}\textgreater{}EX: fills immediats (FILLED/PARTIAL)}
\NormalTok{        EX{-}\textgreater{}\textgreater{}H: record trade agresseur + contreparties}
\NormalTok{    else Ne croise pas}
\NormalTok{        OB{-}{-}\textgreater{}\textgreater{}EX: ACCEPTED + resting order}
\NormalTok{        EX{-}{-}\textgreater{}\textgreater{}API: success, ordre reste dans le carnet}
\NormalTok{    end}

\NormalTok{    TB{-}\textgreater{}\textgreater{}API: POST /trade/EURGBP \{side:buy, qty\} (market)}
\NormalTok{    API{-}\textgreater{}\textgreater{}EX: try\_trade(trader\_id, side, qty, limit\_price=None)}
\NormalTok{    EX{-}\textgreater{}\textgreater{}NPC: update(current\_price)}
\NormalTok{    EX{-}\textgreater{}\textgreater{}OB: execute\_market(...)}
\NormalTok{    OB{-}{-}\textgreater{}\textgreater{}EX: fills par priorite prix/temps}
\NormalTok{    EX{-}\textgreater{}\textgreater{}H: record BUY (Team B) + SELL passif (Team A)}
\NormalTok{    EX{-}{-}\textgreater{}\textgreater{}API: success + average executed price}
\end{Highlighting}
\end{Shaded}

\subsection{Message cle pour la
reunion}\label{message-cle-pour-la-reunion}

\begin{itemize}
\tightlist
\item
  Un \textbf{limit order} n'est pas forcement un trade immediat.
\item
  Il devient trade soit:

  \begin{itemize}
  \tightlist
  \item
    tout de suite s'il croise le carnet,
  \item
    plus tard quand un market order vient le taper.
  \end{itemize}
\item
  Le matching se fait par priorite \textbf{meilleur prix puis ordre
  d'arrivee}.
\item
  Dans l'etat actuel, les deux jambes (buyer et seller) sont bien mises
  a jour en position et historique.
\end{itemize}

\end{document}
